\Conclusion

В ходе выполнения творческой работы разработано автоматизированное рабочее место диспетчера по расписанию. Поставленная цель достигнута, все задачи решены.

Проведён анализ предметной области, спроектирована структура реляционной базы данных, разработана OpenAPI-спецификация REST API. Реализован сервер на Go с автогенерацией кода по спецификации. Разработан десктопный клиент на C++/Qt6 с автогенерацией API-клиента. Реализован импорт расписания из XLSX-файлов с транзакционным механизмом staging-таблиц и экспорт в формат iCalendar. Выполнена контейнеризация приложения средствами Docker~\cite{docker:doc}.

Проект имеет перспективы развития. В будущем возможно добавление автоматического составления расписания в зависимости от заданных параметров: рабочие дни преподавателей, учебное время групп, учёт санитарно-эпидемиологических правил и норм. Также возможна реализация формирования печатного варианта текущей схемы расписания.

Полученные навыки проектирования архитектуры клиент-серверных приложений, работы с OpenAPI, Go, C++/Qt, реляционными базами данных и Docker имеют практическую ценность и могут быть применены при разработке более сложных программных систем.
